\section{Weighted Gaussian arrays and a periodic-memory echo}
\label{sec:weighted-echo}

The extension here concerns the interaction-picture echo, not only returns
of the propagator. These are project proofs; no historical-priority or full
analytic Lean-formalization claim is made.

For a real coefficient array $f=(f_j)_{j=0}^{K-1}$, define
\[
 X_K(f)=\frac1{\sqrt{2K}}\sum_{j=0}^{K-1}f_jB_j,
 \qquad Q_K(f,g)=\frac1K\sum_{j=0}^{K-1}f_jg_j.
\]
For each fixed $r$ and uniformly bounded real arrays $f_1,\ldots,f_{2r}$,
\[
 \begin{split}
 \tau_-\bigl(X_K(f_1)\cdots X_K(f_{2r})\bigr)
 &=\sum_{\pi}\prod_{\{a,b\}\in\pi}Q_K(f_a,f_b)\\
 &\quad+O_r(K^{-1})+O_r(K^r/N),
 \end{split}\tag{W1}
\]
where the sum is over all pairings and the implicit constants may depend
on the common coefficient bound. Odd ordered moments tend to zero.
The arrays may vary with $K$; convergence of their Gram matrices is enough
for a joint centered Gaussian tracial limit.

To prove (W1), the characteristic-zero lifting and bounded-component count
from the preceding argument are unchanged by bounded weights. Nonbalanced
words remain lower order. For each balanced pairing, the coefficient of a
pair $\{a,b\}$ is $f_{a,j}f_{b,j}$. Removing its crossing phase incurs at
most the common weight bound times $9/K$ per crossing, by the summed phase
estimate. The remaining label sums factor into the empirical covariances.
Thus independence of orbit operators or crossing phases is not assumed.

For $p=\max_j|f_j|$, $q=\max_j|g_j|$ and
$c_N=\sqrt{2N/(N-1)}$, disjoint signed-pair Weyl supports also give
\[
 \|[X_K(f),X_K(g)]\|_{2,-}
 \leq\frac{6c_Npq}{\sqrt K},\qquad K\geq11.\tag{W2}
\]
Indeed the full normalized squared norm is
\[
 \frac1{K^2}\sum_{t<s}|f_tg_s-f_sg_t|^2a_{s-t}^2,
\]
which is at most $36p^2q^2/K$. Restriction to the odd subspace gives
(W2). The statement is not an operator-norm almost-commutativity theorem.
The Duhamel estimate and (W2) compare finite products of weighted
exponentials with the exponential of their sum. Applying (W1) to that sum
proves the corresponding Gaussian characteristic-function limits.

\subsection{The echo at every bounded recurrence time}
Let
\[
 \varepsilon=\frac{\tau}{\sqrt{2K}},\qquad
 E_K(t)=U_\varepsilon^tU_0^{-t},\qquad t\in\mathbb Z.
\]
Write $t=mK+s$, $0\leq s<K$, and set
\[
 w_{t,K}(j)=m+\mathbf1_{j<s},\qquad
 c_j=\frac{K-1-2j}{2K}.
\]
The exact identity
\[
 c_j+\frac tK-c_{(j-t)\bmod K}=w_{t,K}(j)\tag{W3}
\]
turns the finite Floquet normal form into a single weighted exponential:
\[
 \|E_K(t)-e^{i\tau X_K(w_{t,K})}\|_{2,-}
 \leq\frac{(11C+2)\tau^2}{\sqrt K},
 \qquad |t|\leq CK,\quad K\geq11.\tag{W4}
\]
For clarity, first telescope the one-step normal form. This approximates
the echo by
\[
 e^{i\tau X_K(c)}
 e^{i\tau(t/K)X_K(1)}
 e^{-i\tau X_K(c_{\mathrm{shift}})}
\]
with error at most $6|t/K|\tau^2/\sqrt K$. Apply (W2) and the two-factor
Duhamel estimate twice. Since $\|c\|_\infty$ and
$\|c_{\mathrm{shift}}\|_\infty$ are at most $1/2$, the extra error is at
most $c_N(3|t/K|+3/4)\tau^2/\sqrt K$. Identity (W3) and
$c_N\leq\sqrt{7/3}$ give (W4). Negative $t$ follows by unitary adjoints
in the telescoping estimate; the division $t=mK+s$ still uses the floor.

For real $\alpha=m+a$, $m\in\mathbb Z$, $0\leq a<1$, define
\[
 Y(\alpha)=m\mathcal B(1)+\mathcal B(a),\tag{W5}
\]
where $\mathcal B$ is one standard real Brownian motion on $[0,1]$.
For $\beta=n+b$ its covariance is
\[
 Q(\alpha,\beta)=mn+mb+na+\min(a,b).\tag{W6}
\]
The same Brownian interval is repeated at every cycle. It is not a sequence
of independent increments across successive cycles.

Fix finitely many real times $\alpha_l$, integers $r_l$, and sequences
$t_{l,K}/K\to\alpha_l$. The finite-dimensional tracial echo limit is
\[
 \tau_-\!\left(\prod_l E_K(t_{l,K})^{r_l}\right)
 \longrightarrow
 \exp\!\left[-\frac{\tau^2}{2}
     \sum_{l,j}r_lr_jQ(\alpha_l,\alpha_j)\right].\tag{W7}
\]
To see this, use (W4) to replace each factor, then (W2) to combine the
weighted exponentials. The step arrays are uniformly bounded and their
Gram matrices converge to (W6), including at integer limiting times.
The weighted Gaussian theorem completes the argument.

For a single time $t_K/K\to\alpha=m+a$, all integer moments of the unitary
echo converge to those of a wrapped normal law with angular variance
\[
 \tau^2\bigl[m^2+(2m+1)a\bigr].\tag{W8}
\]
In particular,
\[
 \tau_-(E_K(t_K))\longrightarrow
 e^{-\tau^2[m^2+(2m+1)a]/2}.
\]
Within the first cycle the variance is $\tau^2\alpha$, whereas at integer
returns it is $\tau^2m^2$. At $\alpha=3/2$ it equals $5\tau^2/2$, not
$3\tau^2/2$ as independent-cycle Brownian noise would predict.

The word Brownian in this statement specifies a limiting tracial Gaussian
covariance. It does not assert a sample-path invariance principle, a joint
spectral distribution for noncommuting finite matrices, or independent
eigenvalues. The time window is still $O(\log N)$ and the classical kick
vanishes. Positive microscopic BGS at fixed classical strength remains a
separate problem.
